%%==========================================================================
%% This is the very first version of a Streetwise LaTeX template. By
%% chosing the document class, it is easily convertible to a two-column IEEE
%% paper format, provided any conflicting packages and (new)commmands are
%% removed.
%%
%% Erwin de Gelder, January 6, 2018
%%==========================================================================

%%==========================================================================
%% Title:
%% Author(s):
%% To be published in:
%% Date of submission:
%% Date of submission R1:
%% Date of submission R2:
%% Date of final submission:
%%==========================================================================


%% Generic
\documentclass[10pt,final,a4paper,oneside,onecolumn]{article}

%% IEEE
%% Document class options: replace "draftcls" by "final" for final document.
%% All other options may just as well be omitted because they are the default values.
%\documentclass[10pt,final,journal,letterpaper,twoside,twocolumn]{IEEEtran}


%%==========================================================================
%% Document automation
%%==========================================================================

\def\reptitle{Similarity measures for traffic scenarios}
\def\repauthor{Erwin de Gelder, etc.}


%%==========================================================================
%% Packages
%%==========================================================================

\usepackage[a4paper,left=3.5cm,right=3.5cm,top=3cm,bottom=3cm]{geometry} %% change page layout; remove for IEEE paper format
\usepackage[T1]{fontenc}                        %% output font encoding for international characters (e.g., accented)
\usepackage[cmex10]{amsmath}                    %% math typesetting; consider using the [cmex10] option
\usepackage{amssymb}                            %% special (symbol) fonts for math typesetting
\usepackage{amsthm}                             %% theorem styles
\usepackage{dsfont}                             %% double stroke roman fonts: the real numbers R: $\mathds{R}$
\usepackage{mathrsfs}                           %% formal script fonts: the Laplace transform L: $\mathscr{L}$
\usepackage[pdftex]{graphicx}                   %% graphics control; use dvips for TeXify; use pdftex for PDFTeXify
\usepackage{array}                              %% array functionality (array, tabular)
\usepackage{upgreek}                            %% upright Greek letters; add the prefix 'up', e.g. \upphi


\usepackage[utf8]{inputenc}    % utf8 support
\usepackage[style=ieee,doi=false,isbn=false,url=false,date=year,minbibnames=15,maxbibnames=15,backend=biber]{biblatex}
%\renewcommand*{\bibfont}{\footnotesize}
\setlength{\biblabelsep}{\labelsep}
\bibliography{../bib}

%\usepackage[round]{natbib}                     %% Author-year citations (remove package cite)
\usepackage{stfloats}                           %% improved handling of floats
\usepackage{multirow}                           %% cells spanning multiple rows in tables
%\usepackage{subfigure}                         %% subfigures and corresponding captions (for use with IEEEconf.cls)
%\usepackage{subfig}                             %% subfigures (IEEEtran.cls: set caption=false)
\usepackage{fancyhdr}                           %% page headers and footers
\usepackage[official,left]{eurosym}             %% the euro symbol; command: \euro
\usepackage{appendix}                           %% appendix layout
\usepackage{xspace}                             %% add space after macro depending on context
\usepackage{verbatim}                           %% provides the comment environment
\usepackage[dutch,USenglish]{babel}             %% language support
\usepackage{wrapfig}                            %% wrapping text around figures
\usepackage{longtable}                          %% tables spanning multiple pages
%\usepackage[utf8]{inputenc}
\usepackage{pgfplots}                           %% support for TikZ figures (Matlab)
\usepackage[breaklinks=true,hidelinks,          %% implement hyperlinks (dvips yields minor problems with breaklinks;
            bookmarksnumbered=true]{hyperref}   %% IEEEtran: set bookmarks=false)
%\usepackage[hyphenbreaks]{breakurl}            %% allow line breaks in URLs (don't use with PDFTeX)
\usepackage{lmodern} 
\usepackage{etoolbox}							%% Needed for apptocmd later
\usepackage[capitalize]{cleveref}
\crefname{figure}{Figure}{Figures}
\crefname{equation}{}{}
\usepackage{units}
\usepackage{subcaption}

%%==========================================================================
%% Fancy headers and footers
%%==========================================================================

\pagestyle{fancy}                                       %% set page style
\fancyhf{}                                              %% clear all header & footer fields
\fancyhead[L]{\includegraphics[width=15mm]{Streetwise}} %% define headers (LE: left field/even pages, etc.)
\fancyhead[R]{\small\emph{\reptitle}}                   %% similar
\fancyfoot[C]{\thepage}                                 %% define footer
\setlength{\headheight}{18pt}                           %% increase head height to accommodate an oversized picture in the header
\renewcommand*{\headrulewidth}{0.25pt}                  %% header line width
\renewcommand*{\footrulewidth}{0pt}                     %% footer line width
%% Redefine the default "plain" page style (automatically activated by \maketitle, \section, ...)
\fancypagestyle{plain}{\fancyhf{}
                       \fancyhead[C]{\includegraphics[width=30mm]{Streetwise}}
                       \renewcommand{\headrulewidth}{0pt}
                       \renewcommand{\footrulewidth}{0pt}
                       \setlength{\headheight}{36pt}}


%%==========================================================================
%% TikZ figures
%%==========================================================================

\newlength\figurewidth
\setlength\figurewidth{0.35\textwidth}              %% set figure height
\newlength\figureheight
\setlength\figureheight{0.3\textwidth}              %% set figure width
\pgfplotsset{every axis/.append style={
    scaled y ticks=false,
    scaled x ticks=false,
    y tick label style={/pgf/number format/fixed},
    x tick label style={/pgf/number format/fixed},
    legend style={font=\small}},
    compat=1.9}                                     %% PGFPlots package options
\usetikzlibrary{shapes.geometric, arrows, arrows.meta}
\usepgfplotslibrary{groupplots}
%\usetikzlibrary{external}                           %% Create pdf figures from TikZ. Use PDFTeXify ...
%\tikzexternalize[prefix=./tikz/]                    %% ... with --tex-option=--shell-escape switch.
%\tikzset{external/force remake}                    %% force pdf figure update


%%==========================================================================
%% User-defined commands
%%==========================================================================

\newcommand*{\mat}[1]{\mathbf{#1}}                              %% matrix/vector notation
\newcommand*{\matsym}[1]{\boldsymbol{#1}}                       %% matrix/vector notation for Greek letters
\newcommand*{\T}{^{\scriptscriptstyle\mathsf{T}}}               %% transpose operator
\newcommand*{\Hr}{^{\scriptscriptstyle\mathsf{H}}}              %% conjugate transpose operator
\newcommand*{\ud}{\mathrm{\,d}}                                 %% differential operator (upright d)
\newcommand*{\defeq}{\mathrel{\mathop:}=}                       %% definition sign :=
\newcommand*{\eqdef}{=\mathrel{\mathop:}}                       %% definition sign =:
\newcommand*{\ip}[2]{\left\langle#1\,{,}\,#2\right\rangle}      %% inner product
\newcommand*{\real}[1]{\mathrm{Re}(#1)}                         %% real part
\newcommand*{\imag}[1]{\mathrm{Im}(#1)}                         %% imaginary part
\newcommand*{\lsup}[1]{{}^{#1}\!}                               %% left superscript
\newcommand*{\hi}[1]{$^\text{#1}$}                              %% superscript in normal text
\newcommand*{\lo}[1]{$_\text{#1}$}                              %% subscript in normal text
\newcommand*{\w}[1]{\mathrm{#1}}                                %% multiple character super-/subscript in math mode
\newcommand*{\Ln}{{\ensuremath{\mathcal{L}}}\xspace}            %% L norm
\newcommand*{\Lone}{{\ensuremath{\mathcal{L}_1}}\xspace}        %% L_1 norm
\newcommand*{\Ltwo}{{\ensuremath{\mathcal{L}_2}}\xspace}        %% L_2 norm
\newcommand*{\Linf}{{\ensuremath{\mathcal{L}_\infty}}\xspace}   %% L_inf norm
\newcommand*{\Lp}{{\ensuremath{\mathcal{L}_p}}\xspace}          %% L_p norm
\newcommand*{\Htwo}{{\ensuremath{\mathcal{H}_2}}\xspace}        %% H_2 norm
\newcommand*{\Hinf}{{\ensuremath{\mathcal{H}_\infty}}\xspace}   %% H_inf norm
\newcommand*{\Hp}{{\ensuremath{\mathcal{H}_p}}\xspace}          %% H_p norm
\newcommand*{\capskip}{\vspace{-12pt}}                          %% caption skip for figures with subfloats
\newcommand*{\etal}{et al.}                                     %% may be required for Natbib bibliography styles
\renewcommand*{\qedsymbol}{$\blacksquare$}                      %% redefine the end-of-proof symbol
\renewcommand*{\labelitemi}{$\bullet$}                          %% first level item list bullet
\renewcommand*{\labelitemii}{$-$}                               %% second level item list bullet
%\renewcommand*{\theenumi}{\textit{\roman{enumi}}}               %% first level enumerator
\renewcommand*{\labelenumi}{\theenumi.}
\renewcommand*{\theenumii}{\textit{\alph{enumii}}}              %% second level enumerator
\renewcommand*{\labelenumii}{\theenumii.}
\DeclareMathOperator{\tr}{tr}                                   %% trace of a matrix
\DeclareMathOperator{\sgn}{sgn}                                 %% signum function
\DeclareMathOperator{\atan}{atan}                               %% arc tangent

\newcommand{\deuclid}{d_{\textup{Euclidian}}}
\newcommand{\ddtw}{d_{\textup{DTW}}}
\newcommand{\imagi}{\textbf{j}}
\newcommand{\profile}[1]{\textbf{v}_{#1}}
\newcommand{\barprofile}[1]{\bar{\textbf{v}}_{#1}}
\newcommand{\parvstart}{v_{\textup{start}}}
\newcommand{\parvend}{v_{\textup{end}}}
\newcommand{\pardv}{v_{\Delta}}
\newcommand{\pardt}{t_{\textup{duration}}}


%%==========================================================================
%% User-defined environments
%%==========================================================================

\theoremstyle{plain}\newtheorem{definition}{Definition}[section]    %% definition
                    \newtheorem{theorem}{Theorem}[section]          %% theorem
                    \newtheorem{lemma}[theorem]{Lemma}              %% lemma
                    \newtheorem{corollary}[theorem]{Corollary}      %% corollary
                    \newtheorem{assumption}{Assumption}[section]    %% assumption
                    \newtheorem{condition}{Condition}[section]      %% condition
\theoremstyle{definition}\newtheorem{example}{Example}[section]     %% examples
\theoremstyle{remark}\newtheorem{remarkenv}{Remark}[section]        %% remarks
\newenvironment{remark}{\begin{remarkenv}}%
                       {\hfill$\blacklozenge$\end{remarkenv}}       %% end remark with a lozenge
\newenvironment{reviewer}{\itshape}{\upshape}                       %% environment for reviewer's comments


%%==========================================================================
%% Miscellaneous
%%==========================================================================

\graphicspath{{./../}{./figures/}{./more_figures/}} %% (graphicx) directory path for figures
%\setlength{\parindent}{0pt}                        %% no paragraph indentation
%\setlength{\parskip}{2ex}                          %% create empty line between paragraphs
%\interdisplaylinepenalty=2500                      %% (amsmath) allow for page breaks within multiline equations
%\numberwithin{equation}{section}                   %% (amsmath) include section number in equation numbering
%\numberwithin{figure}{section}                     %% (amsmath) include section number in figure numbering
%\numberwithin{table}{section}                      %% (amsmath) include section number in table numbering
\addtolength{\arraycolsep}{-0.5mm}                  %% squeeze matrix columns a little
\fnbelowfloat                                       %% (stfloats) put footnote below a float at the page bottom
\urlstyle{same}                                     %% (hyperref) use current font for URLs
\hypersetup{pdftitle={\reptitle},
            pdfauthor={\repauthor}}                 %% (hyperref) pdf properties title and author
%\raggedbottom                                      %% don't add inter-paragraph spacing to achieve \textheight
%\setlength\subfigcapskip{-3pt}                     %% (subfigure) distance between subfloat and subcaption
%\setlength\subfigbottomskip{-3pt}                  %% (subfigure) distance between subcaption and caption
%\renewcommand*{\subcapsize}{\small}                %% (subfigure) subcaption font size
\renewcommand*{\thesubfigure}{(\alph{subfigure})}   %% (subfig) implement 1(a) instead of 1a ...as subfigure reference
\captionsetup[subfloat]{labelformat=simple}        %% ... as subfigure reference
%\captionsetup[subfloat]{
%    farskip=0pt,
%    nearskip=8pt,
%    captionskip=1.5pt,
%    labelfont={small,bf},
%    textfont=small}                                %% (subfig) subfloat caption format
%\captionsetup[figure]{
%    labelfont={small,bf},
%    textfont=small}                                %% (subfig/caption) figure caption format
%\captionsetup[table]{
%    aboveskip=2pt,
%    labelfont={small,bf},
%    textfont={small,sc}}                           %% (subfig/caption) table caption format
\apptocmd{\thebibliography}{\raggedright}{}{}		%% Suppress badness warnings in bibliography


% Table stuff
\usepackage{booktabs}
\usepackage{tabularx}
\setlength{\heavyrulewidth}{0.1em}
\newcommand{\otoprule}{\midrule[\heavyrulewidth]}

%%==========================================================================
%% Begin document
%%==========================================================================

\begin{document}

\selectlanguage{USenglish}

\title{\textbf{\reptitle}}
\author{\repauthor}
\date{\today}
\maketitle

\tableofcontents

\newpage

\section{Introduction}
\label{sec:introduction}
% Scenario-based approach for assessment of AVs -> this scenarios are important
An important aspect in the development of automated vehicles (AVs) is the assessment of quality and performance aspects of the AVs, such as safety, comfort, and efficiency \cite{bengler2014threedecades, stellet2015taxonomy, putz2017pegasus, roesener2016scenariobased, wachenfeld2016release, kompass2015sicherheitsveranderung}. 
For legal and public acceptance, it is important that there is a clear definition of system performance and that there are quantitative measures for the system quality. 
The more traditional methods \cite{response2006code, ISO26262}, used for evaluation of driver assistance systems, are no longer valid for the assessment of quality and performance aspects of an AV \cite{wachenfeld2016release}. 
Therefore, a scenario-based approach is proposed \cite{roesener2016scenariobased, putz2017pegasus, kompass2015sicherheitsveranderung}. 
For the scenario-based assessment, proper specification of scenarios is crucial since they are directly reflected in test cases used for scenario-based assessment \cite{stellet2015taxonomy}.

% Purpose of document
To structure the variety of scenarios, the notion of scenarios classes is introduced \cite{degelder2018ontology}. Whereas a scenario is a quantitative description, a scenario class is a qualitative description \cite{degelder2018ontology}. This document describes different measures to quantify the (dis)similarity between two different scenarios that belong to the same scenario class. An appropriate (dis)similarity measure for scenarios would be useful for at least two reasons.

% Use of a measure for completeness
The first use of a (dis)similarity measure is for quantifying the so-called \emph{completeness}. In order to draw conclusions on how an automated vehicle would perform in real-world traffic, it is necessary to know how representative the scenario catalog, which is used for the scenario-based assessment of the automated vehicle, is. Therefore it is important to quantify how complete the scenario catalog is \cite{geyer2014, alvarez2017prospective, stellet2015taxonomy}. A (dis)similarity measure of scenarios is useful for quantifying the completeness. For example, if newly collected scenarios are not similar to earlier collected scenarios, it is an indicated that the set of collected scenarios is not \emph{complete}. To achieve a complete set of scenarios, this process could be repeated until the newly collected scenarios are - to a certain extent - similar to the early collected scenarios. This process of achieving a complete set of scenarios is visualized in \cref{fig:completeness collection}.

\begin{figure}[b]
	\centering
	\tikzstyle{startstop}=[rectangle, rounded corners, minimum width=4em, minimum height=4em, text centered, draw=black, fill=red!30, text width=4em]
	\tikzstyle{process} = [rectangle, minimum width=9em, text width=8.5em, minimum height=4em, text centered, draw=black, fill=orange!30, node distance=10em]
	\tikzstyle{decision} = [diamond, minimum width=4em, text width=4em, minimum height=4em, text centered, draw=black, fill=green!30, node distance=10em]
	\tikzstyle{arrow} = [-{Triangle[width=2mm, length=3mm, round]}, line width=0.6mm] % rounded corners=0.4em
	\begin{tikzpicture}
		% Draw the blocks
		\node[startstop](start){Start};
		\node[process, right of=start, node distance=7em, text width=5em, minimum width=5.5em](collect){Collect new scenarios};
		\node[process, right of=collect, node distance=9em](compare){Compare new scenarios with scenario database};
		\node[decision, right of=compare](similar){New scenarios similar?};
		\node[startstop, right of=similar, node distance=9em](stop){Scenario database complete};
		\node[process, below of=compare, node distance=7em](add){Add collected scenarios to scenario database};
		
		% Arrows
		\draw[arrow] (start) -- (collect);
		\draw[arrow] (collect) -- (compare);
		\draw[arrow] (compare) -- (similar);
		\draw[arrow] (similar) -- node[anchor=south]{Yes} (stop);
		\draw[arrow] (similar) |- node[anchor=north]{No} (add);
		\draw[arrow] (add) -| (collect);
	\end{tikzpicture}
	\caption{Schematic process of collecting scenarios of a specific \emph{scenario class} until a certain completeness is reached.}
	\label{fig:completeness collection}
\end{figure}

% Use of a measure for scenario generation
The second use of a (dis)similarity measure is for the generation of test scenarios. Using collected real-world scenarios, one might want to generate an additional set of scenarios to be used for the scenario-based assessment in order to emphasize safety-critical or performance-critical scenarios (cf. \cite{zofka2015datadrivetrafficscenarios, deGelder2017assessment}). Although the generated scenarios are different from the collected scenarios, it is still desirable that the generated scenarios are reflecting what might be encountered in real-life. For this purpose, a (dis)similarity measure can be used. For example, if the generated scenarios is similar to one or several encountered scenarios, it is more conceivable that this scenario is a realistic scenario compared to a scenario that does not show high similarities with any of the collected scenarios. 


\section{Dissimilarity measures}
\label{sec:dissimilarity}

In this section, several dissimilarity measures are presented. They can be divided into four different categories \cite{serra2014empirical}: lock-step measures, feature-based measures, model-based measures, and elastic measures. Before discussing these different categories, some necessary notation details are described.

\subsection{Notation and example}
\label{sec:dissimilarity notation}

It is assumed that the scenario is described by means of a time series. Therefore, the dissimilarity measures compare the time series $\textbf{x}$ and $\textbf{y}$. The i-th element of $\textbf{x}$ and $\textbf{y}$ are denoted by $x_i$ and $y_i$, respectively. For the sake of simplicity, $x_i$ and $y_i$ are scalars. The length of the time series $\textbf{x}$ and $\textbf{y}$ is $N$ and $M$, respectively. Note that the time series do not necessarily have the same length, i.e., $N$ might be different than $M$. It is assumed, however, that the sample time is constant and similar for both time series.

To illustrate the dissimilarity measures, four different velocity profiles are considered, see \cref{fig:example time series}. All four velocity profiles describe the velocity of a vehicle while braking. The average deceleration is $\unit[3.19]{m/s^2}$, $\unit[1.72]{m/s^2}$, $\unit[1.63]{m/s^2}$, and $\unit[1.64]{m/s^2}$ respectively. We will refer to the the profiles with $\profile{1}$, $\profile{2}$, $\profile{3}$, and $\profile{4}$, respectively. 

The time series are chosen for the sake of illustrating the features of the dissimilarity measures. Time series $\profile{1}$ has the same start and end speed as the time series $\profile{2}$ and $\profile{3}$, but the end speed is earlier achieved, i.e., the deceleration is stronger. The time series $\profile{2}$ and $\profile{3}$ are reasonably similar, so a low dissimilarity score is to be expected when comparing $\profile{2}$ and $\profile{3}$. The fourth time series, i.e., $\profile{4}$, is chosen because its average deceleration almost equals the average deceleration of $\profile{3}$, but the number of samples of $\profile{4}$ is much less.

\begin{figure}
	\centering
	\setlength\figureheight{0.4\linewidth}
	\setlength\figurewidth{0.7\linewidth}
	% This file was created by matplotlib2tikz v0.6.14.
\begin{tikzpicture}

\definecolor{color3}{rgb}{0.83921568627451,0.152941176470588,0.156862745098039}
\definecolor{color2}{rgb}{0.172549019607843,0.627450980392157,0.172549019607843}
\definecolor{color1}{rgb}{1,0.498039215686275,0.0549019607843137}
\definecolor{color0}{rgb}{0.12156862745098,0.466666666666667,0.705882352941177}

\begin{axis}[
xlabel={Time [s]},
ylabel={Velocity [km/h]},
xmin=-0.405, xmax=8.505,
ymin=-0.128527297559858, ymax=50.9958346332171,
width=\figurewidth,
height=\figureheight,
tick align=outside,
tick pos=left,
xmajorgrids,
x grid style={white!69.019607843137251!black},
ymajorgrids,
y grid style={white!69.019607843137251!black},
legend cell align={left},
legend entries={{Profile 1},{Profile 2},{Profile 3},{Profile 4}},
legend style={draw=white!80.0!black}
]
\addlegendimage{no markers, color0}
\addlegendimage{no markers, color1}
\addlegendimage{no markers, color2}
\addlegendimage{no markers, color3}
\addplot [semithick, color0, mark=*, mark size=1, mark options={solid}, only marks]
table {%
0 48.672
0.1 48.6359918994937
0.2 47.9159878492406
0.3 47.375962482237
0.4 46.4759661771686
0.5 45.5398532671499
0.6 44.567824915378
0.7 43.4157646613031
0.8 42.29969559451
0.9 41.2916970126468
1 39.9956646106218
1.1 38.8618378493623
1.2 37.5478254838316
1.3 36.4676637612536
1.4 35.1716329227879
1.5 34.0015198318347
1.6 32.8674551343703
1.7 31.4097739319247
1.8 29.9876879204179
1.9 28.7453026814198
2 27.467341656224
2.1 26.2253887315345
2.2 24.9472728019208
2.3 23.471640170536
2.4 22.3550563713891
2.5 20.6271142735906
2.6 19.295333982059
2.7 17.8910381809574
2.8 16.5229428943627
2.9 15.1007890808964
3 13.7507984249087
3.1 12.2935431385527
3.2 10.9075264803432
3.3 9.8270245995681
3.4 8.78306794438217
3.5 7.84696789602428
3.6 6.87493954425236
3.7 5.81345310924209
3.8 4.78727862860074
3.9 3.83293169244461
4 2.98698956843009
4.1 2.19530733565728
};
\addplot [semithick, color1, mark=*, mark size=1, mark options={solid}, only marks]
table {%
0 47.088
0.1 46.9079863570422
0.2 46.7999676690327
0.3 46.5479802461338
0.4 46.3679702254707
0.5 46.1159598531504
0.6 45.9358601596821
0.7 45.4498142231174
0.8 45.179940454173
0.9 44.729704544694
1 44.3158188044651
1.1 43.8836856435103
1.2 43.3796156522959
1.3 42.767789955294
1.4 42.0838333001078
1.5 41.5075281810371
1.6 40.8955037374065
1.7 40.2294110079255
1.8 39.5993835088374
1.9 38.8795685414562
2 37.9797001391539
2.1 37.3852946164394
2.2 36.68367427438
2.3 35.9631533633361
2.4 35.3150331264124
2.5 34.5051798318338
2.6 33.6233394534536
2.7 32.9032192325912
2.8 32.111495850071
2.9 31.4444417183995
3 30.4553562229219
3.1 29.9144526054014
3.2 29.2667486138953
3.3 28.3131804140041
3.4 27.5397595428659
3.5 26.8006419351036
3.6 26.2428597471555
3.7 25.5592204405501
3.8 24.8025352295064
3.9 24.0298637835918
4 23.3986030116011
4.1 22.6793546660349
4.2 21.9596708636399
4.3 21.4004881897229
4.4 20.7706317250088
4.5 20.2666044390949
4.6 19.6907574960284
4.7 19.1152763547719
4.8 18.5036320664783
4.9 17.9985009799961
5 17.5308574022717
5.1 17.0268369378363
5.2 16.4870142531641
5.3 15.8930926011825
5.4 15.299260445287
5.5 14.7406843782999
5.6 14.1466629901203
5.7 13.5526391149456
5.8 12.958617726766
5.9 12.453991216557
6 11.8423685296134
6.1 11.1595853109241
6.2 10.4393688000234
6.3 9.80807570341702
6.4 9.21361447195764
6.5 8.53068047016063
6.6 7.84599002842424
6.7 7.10886806083318
6.8 6.44308081477942
6.9 5.81190431375665
7 5.14634863363259
7.1 4.56960679779292
7.2 3.92181861892924
7.3 3.2390168543433
7.4 2.51875579065799
};
\addplot [semithick, color2, mark=*, mark size=1, mark options={solid}, only marks]
table {%
0 48.672
0.1 48.608995283563
0.2 48.455994244377
0.3 48.1499559446157
0.4 48.0239525338739
0.5 47.7179552340439
0.6 47.4119101119168
0.7 47.0879575785089
0.8 46.6558044509344
0.9 46.1519450018259
1 45.6478472272915
1.1 45.2157983400267
1.2 44.8197433393213
1.3 44.3876808084788
1.4 43.9557843417831
1.5 43.6319079801829
1.6 43.0918501391549
1.7 42.6235759313898
1.8 42.1916050647879
1.9 41.651644213726
2 41.0034997582104
2.1 40.5357376572865
2.2 39.9237254327615
2.3 39.2933547307219
2.4 38.6992524722155
2.5 38.0150641691167
2.6 37.4394314723628
2.7 36.773082805036
2.8 36.0714011250792
2.9 35.2613784931797
3 34.5953583483032
3.1 33.785335716398
3.2 33.1913130912778
3.3 32.5071905178296
3.4 32.2546217634148
3.5 31.2471259893812
3.6 30.4914375832673
3.7 29.8430756895968
3.8 29.2670512366027
3.9 28.6192701017491
4 28.1511270552457
4.1 27.5027206047803
4.2 26.9273441432204
4.3 26.4231946391366
4.4 25.8835881532068
4.5 25.4154386491253
4.6 24.8389962378834
4.7 24.2806803505508
4.8 23.7949509844374
4.9 23.2190825110755
5 22.6078440297254
5.1 22.1208868199009
5.2 21.6710147129051
5.3 21.1306604188455
5.4 20.6268052463816
5.5 20.033131327285
5.6 19.4392928330335
5.7 18.8985612363945
5.8 18.3587172066642
5.9 17.7468831849665
6 17.1702956961133
6.1 16.5410378872341
6.2 15.8750226003189
6.3 15.3352051477411
6.4 14.6869920454748
6.5 14.1281662850803
6.6 13.4985585362236
6.7 12.9769551057457
6.8 12.3463037088183
6.9 11.6991396409641
7 11.0151274105371
7.1 10.3302289560398
7.2 9.6455292745252
7.3 8.87149566453292
7.4 8.0259234991623
7.5 7.19859733338512
7.6 6.37057898663801
7.7 5.52335924441851
7.8 4.71343313430595
7.9 3.90339952431368
8 3.20136591431246
8.1 2.37502992884368
};
\addplot [semithick, color3, mark=*, mark size=1, mark options={solid}, only marks]
table {%
0 30.492
0.1 30.3839935337426
0.2 30.0959935337426
0.3 29.7719418036832
0.4 29.411922404911
0.5 28.9439030061386
0.6 28.4038254110494
0.7 27.719954736198
0.8 27.0358189447921
0.9 26.3157090184159
1 25.6315473619803
1.1 24.8756443558417
1.2 24.1555732312219
1.3 23.4355796932674
1.4 22.715502103092
1.5 21.9058545092079
1.6 21.1317930715936
1.7 20.4113233878768
1.8 19.6554038252768
1.9 18.9894343445907
2 18.3954052464323
2.1 17.6397227325662
2.2 17.0278549000257
};
\end{axis}

\end{tikzpicture}
	\caption{Four examples of time series which will be used to illustrate the dissimilarity measures. The velocity is sampled with $\unit[10]{Hz}$.}
	\label{fig:example time series}
\end{figure}

\subsection{Lock-step measures}
\label{sec:lock step}

Lock-step measures estimate the dissimilarity between two time series by comparing each time instant. This, however, does only work if the two time series have equal length. In case that the time series $\textbf{x}$ and $\textbf{y}$ do not have the same length, one of the time series can be re-sampled. This works well for a number of data sources \cite{keogh2003survey}. Let the time series $\bar{\textbf{y}}$ denote the re-sampled time series $\textbf{y}$, such that $\bar{\textbf{y}}$ contains $N$ samples, e.g., using linear interpolation. As a consequence, the elements $\bar{y}[i]$ of $\bar{\textbf{y}}$ are not sampled with the same frequency as the time series $\textbf{x}$ and $\textbf{y}$. 

A straightforward way to estimate the dissimilarity between two time series is to use an $L_p$ norm. The $L_p$ norm is as follows:
\begin{equation} \label{eq:lp norm}
	d_{L_p}(\textbf{x}, \bar{\textbf{y}}) = \left( \sum_{i=0}^{N-1} |x[i] - \bar{y}[i]|^p \right)^{\frac{1}{p}},
\end{equation}
where $p \geq 0$. Not that with $p=2$, we obtain the Euclidean distance:
\begin{equation} \label{eq:euclidian}
	\deuclid(\textbf{x}, \bar{\textbf{y}}) = \sqrt{\sum_{i=0}^{N-1} |x[i] - \bar{y}[i]|^p}.
\end{equation}
The Euclidean distance is often applied due to its computational simplicity \cite{xing2010survey}. Even if there might be better dissimilarity measures for certain applications, the Euclidean distance serves as a good benchmark \cite{keogh2003survey}.

For our purposes, the Euclidean distances is not suitable. Because, in general, we have $N \ne M$, we need to re-sample one time series. Because of this, a strong braking action might look `similar' to very smooth braking, see \cref{fig:resampled profiles}. In \cref{fig:resampled profiles}, $\textbf{v}_1$ and $\textbf{v}_3$ are re-sample, such that they have the same amount of samples as $\profile{2}$, resulting in $\barprofile{1}$ and $\barprofile{3}$, respectively. This results in $\deuclid(\profile{2}, \barprofile{1})=\unit[4.7]{m/s}$ and $\deuclid(\profile{2}, \barprofile{3})=\unit[4.2]{m/s}$. So even though $\profile{3}$ looks much more similar to $\profile{2}$ than $\profile{1}$, the Euclidean distance \cref{eq:euclidian} gives not a significant difference between $\deuclid(\profile{2}, \barprofile{3})$ and $\deuclid(\profile{2}, \barprofile{1})$.

\begin{figure}
	\centering
	\setlength\figureheight{0.4\linewidth}
	\setlength\figurewidth{0.46\linewidth}
	\begin{minipage}{0.46\linewidth}
		\centering
		% This file was created by matplotlib2tikz v0.6.14.
\begin{tikzpicture}

\definecolor{color1}{rgb}{1,0.498039215686275,0.0549019607843137}
\definecolor{color0}{rgb}{0.12156862745098,0.466666666666667,0.705882352941177}

\begin{axis}[
xlabel={Sample},
ylabel={Velocity [km/h]},
xmin=-3.7, xmax=77.7,
ymin=-0.128527297559858, ymax=50.9958346332171,
width=\figurewidth,
height=\figureheight,
tick align=outside,
tick pos=left,
xmajorgrids,
x grid style={white!69.019607843137251!black},
ymajorgrids,
y grid style={white!69.019607843137251!black}
]
\addplot [semithick, color0, mark=*, mark size=1, mark options={solid}, only marks, forget plot]
table {%
0 47.088
1 46.9079863570422
2 46.7999676690327
3 46.5479802461338
4 46.3679702254707
5 46.1159598531504
6 45.9358601596821
7 45.4498142231174
8 45.179940454173
9 44.729704544694
10 44.3158188044651
11 43.8836856435103
12 43.3796156522959
13 42.767789955294
14 42.0838333001078
15 41.5075281810371
16 40.8955037374065
17 40.2294110079255
18 39.5993835088374
19 38.8795685414562
20 37.9797001391539
21 37.3852946164394
22 36.68367427438
23 35.9631533633361
24 35.3150331264124
25 34.5051798318338
26 33.6233394534536
27 32.9032192325912
28 32.111495850071
29 31.4444417183995
30 30.4553562229219
31 29.9144526054014
32 29.2667486138953
33 28.3131804140041
34 27.5397595428659
35 26.8006419351036
36 26.2428597471555
37 25.5592204405501
38 24.8025352295064
39 24.0298637835918
40 23.3986030116011
41 22.6793546660349
42 21.9596708636399
43 21.4004881897229
44 20.7706317250088
45 20.2666044390949
46 19.6907574960284
47 19.1152763547719
48 18.5036320664783
49 17.9985009799961
50 17.5308574022717
51 17.0268369378363
52 16.4870142531641
53 15.8930926011825
54 15.299260445287
55 14.7406843782999
56 14.1466629901203
57 13.5526391149456
58 12.958617726766
59 12.453991216557
60 11.8423685296134
61 11.1595853109241
62 10.4393688000234
63 9.80807570341702
64 9.21361447195764
65 8.53068047016063
66 7.84599002842424
67 7.10886806083318
68 6.44308081477942
69 5.81190431375665
70 5.14634863363259
71 4.56960679779292
72 3.92181861892924
73 3.2390168543433
74 2.51875579065799
};
\addplot [semithick, color1, mark=*, mark size=1, mark options={solid}, only marks, forget plot]
table {%
0 48.672
1 48.6515629699829
2 48.5386940548649
3 48.130043107424
4 47.7700350473477
5 47.4635341633727
6 47.0110991153174
7 46.5002904016299
8 45.9699591987801
9 45.4347691210124
10 44.8830773537905
11 44.2875940427652
12 43.6337220066686
13 42.9934682576516
14 42.3600236521745
15 41.7820747011208
16 41.1866133043745
17 40.4510273464684
18 39.7505128784575
19 39.1069895815265
20 38.4001578290407
21 37.6543670269827
22 37.0223414025774
23 36.3976080402554
24 35.6620229697749
25 34.9502601758508
26 34.2861419350395
27 33.6337150650895
28 32.9900567232854
29 32.1977097170304
30 31.371339174857
31 30.5642092764342
32 29.7862200438236
33 29.0810824757436
34 28.3653683225778
35 27.6400390920613
36 26.9302809320339
37 26.2253887315345
38 25.4999715822943
39 24.7478629868688
40 23.9103417636504
41 23.1698607653612
42 22.536124014494
43 21.654539304714
44 20.6738154113689
45 19.9072330349248
46 19.1435182197777
47 18.3464854678012
48 17.5582582463803
49 16.7817717323671
50 15.9848306406187
51 15.1776622600027
52 14.4075506359298
53 13.6326425908799
54 12.8055517526778
55 11.993863861102
56 11.2072057577939
57 10.5278906843952
58 9.91463286017149
59 9.31915379434252
60 8.73246794176824
61 8.20116791432186
62 7.66307064028365
63 7.11137887306175
64 6.53067367343822
65 5.92820839951347
66 5.34196753705552
67 4.76148546816409
68 4.21982909899439
69 3.69575188855036
70 3.2156225749205
71 2.75162458030844
72 2.30229142116712
73 2.19530733565728
74 2.19530733565728
};
\end{axis}

\end{tikzpicture}
		\subcaption{$\deuclid(\textbf{v}_2, \bar{\textbf{v}}_1)=\unit[13.3]{m/s}$, $\textbf{v}_2$ in blue and $\bar{\textbf{v}}_1$ in orange.}
	\end{minipage}
	\hspace{0.04\linewidth}
	\begin{minipage}{0.46\linewidth}
		% This file was created by matplotlib2tikz v0.6.14.
\begin{tikzpicture}

\definecolor{color1}{rgb}{1,0.498039215686275,0.0549019607843137}
\definecolor{color0}{rgb}{0.12156862745098,0.466666666666667,0.705882352941177}

\begin{axis}[
xlabel={Sample},
ylabel={Velocity [km/h]},
xmin=-3.7, xmax=77.7,
ymin=0.0601814252858692, ymax=50.9868485035578,
width=\figurewidth,
height=\figureheight,
tick align=outside,
tick pos=left,
xmajorgrids,
x grid style={white!69.019607843137251!black},
ymajorgrids,
y grid style={white!69.019607843137251!black}
]
\addplot [semithick, color0, mark=*, mark size=1, mark options={solid}, only marks, forget plot]
table {%
0 47.088
1 46.9079863570422
2 46.7999676690327
3 46.5479802461338
4 46.3679702254707
5 46.1159598531504
6 45.9358601596821
7 45.4498142231174
8 45.179940454173
9 44.729704544694
10 44.3158188044651
11 43.8836856435103
12 43.3796156522959
13 42.767789955294
14 42.0838333001078
15 41.5075281810371
16 40.8955037374065
17 40.2294110079255
18 39.5993835088374
19 38.8795685414562
20 37.9797001391539
21 37.3852946164394
22 36.68367427438
23 35.9631533633361
24 35.3150331264124
25 34.5051798318338
26 33.6233394534536
27 32.9032192325912
28 32.111495850071
29 31.4444417183995
30 30.4553562229219
31 29.9144526054014
32 29.2667486138953
33 28.3131804140041
34 27.5397595428659
35 26.8006419351036
36 26.2428597471555
37 25.5592204405501
38 24.8025352295064
39 24.0298637835918
40 23.3986030116011
41 22.6793546660349
42 21.9596708636399
43 21.4004881897229
44 20.7706317250088
45 20.2666044390949
46 19.6907574960284
47 19.1152763547719
48 18.5036320664783
49 17.9985009799961
50 17.5308574022717
51 17.0268369378363
52 16.4870142531641
53 15.8930926011825
54 15.299260445287
55 14.7406843782999
56 14.1466629901203
57 13.5526391149456
58 12.958617726766
59 12.453991216557
60 11.8423685296134
61 11.1595853109241
62 10.4393688000234
63 9.80807570341702
64 9.21361447195764
65 8.53068047016063
66 7.84599002842424
67 7.10886806083318
68 6.44308081477942
69 5.81190431375665
70 5.14634863363259
71 4.56960679779292
72 3.92181861892924
73 3.2390168543433
74 2.51875579065799
};
\addplot [semithick, color1, mark=*, mark size=1, mark options={solid}, only marks, forget plot]
table {%
0 48.672
1 48.592454630678
2 48.3898238011854
3 48.1090899735643
4 47.8916293771907
5 47.5525254382995
6 47.2017787388955
7 46.7609227792633
8 46.2200341165703
9 45.6614714914681
10 45.183685772402
11 44.7380017794322
12 44.2592791562179
13 43.8244831141074
14 43.3545809807361
15 42.8007607667604
16 42.3083539476533
17 41.7392054328171
18 41.0385345936437
19 40.5026559154203
20 39.821503156755
21 39.1327865527472
22 38.4403704115835
23 37.7350266409662
24 37.043224156655
25 36.2800091920934
26 35.4146260181337
27 34.6493599816715
28 33.7692810508542
29 33.0988640948658
30 32.4457548748638
31 31.9006367617273
32 30.8999178027883
33 30.1234484003732
34 29.4538699781143
35 28.7593308876634
36 28.2017371143271
37 27.5027206047803
38 26.8728414941303
39 26.3065229665031
40 25.7317558816128
41 25.1661662550748
42 24.5372038663523
43 23.9656126536124
44 23.3591586262176
45 22.6904438245024
46 22.1340478255719
47 21.627202202576
48 21.0353364672983
49 20.4503075947583
50 19.7923859917776
51 19.1616198509756
52 18.5629825152108
53 17.9122437313713
54 17.2637963699814
55 16.5750518228492
56 15.8458432785579
57 15.2300895095358
58 14.535958056179
59 13.8899363260535
60 13.2448055159911
61 12.6019731940591
62 11.8915397692451
63 11.1445351298071
64 10.3857612631612
65 9.62460944722811
66 8.75722915569906
67 7.82468199937867
68 6.90767196831181
69 5.98131586183445
70 5.06367145219246
71 4.16611312755441
72 3.35315696512353
73 2.46436354889437
74 2.37502992884368
};
\end{axis}

\end{tikzpicture}
		\subcaption{$\deuclid(\textbf{v}_2, \bar{\textbf{v}}_3)=\unit[14.7]{m/s}$, $\textbf{v}_2$ in blue and $\bar{\textbf{v}}_3$ in orange.}
	\end{minipage}
	\caption{Comparison of the different profiles, where one of the two profiles is re-sampled.}
	\label{fig:resampled profiles}
\end{figure}


\subsection{Feature-based measures}
\label{sec:feature based}

The Euclidean distance of \cref{eq:euclidian} is not of direct use for our application. An simple extension, however, is useful. This extension is to compute the Euclidean distance not directly using the (re-sampled) time series, but using features extracted from it \cite{serra2014empirical}. Let $\hat{x}_k$ and $\hat{y}_k$ with $k \in \{0, \ldots, \theta-1\}$ denote features extracted from the time series $\textbf{x}$ and $\textbf{y}$, respectively. The dissimilarity can be formulated as follows:

\begin{equation}\label{eq:eucl features}
	d_{\textup{F}}(\textbf{x}, \textbf{y}) = \sqrt{ \sum_{k=0}^{\theta-1} w_k \left( \hat{x}_k - \hat{y}_k \right)^2 }.
\end{equation}
Here, $w_k$ denote weights that can be chosen to emphasize the $k$-th parameter less or more. 


\subsubsection{Fourier coefficients}
\label{sec:fourier coefficient}
One possibility is to use the Fourier coefficients as parameters \cite{agrawal1993similarity, faloutsos1994fast}. The Fourier coefficients are obtained with the Discrete Fourier Transform \cite{oppenheim1999discrete}. The $n$-th Fourier coefficient, i.e.,
\begin{equation}
	X[n] = \sum_{i=0}^{N-1} x[i] e^{-\imagi\frac{2\pi i n}{N}},
\end{equation}
correspond to a frequency of $f_n = n/Nt_s$, where $N$ is the total number of samples and $t_s$ is the sample time. When only the Fourier coefficients corresponding to the lower frequencies are used as parameters, this automatically has the effect of removing rapidly-fluctuating components of the time series \cite{serra2014empirical}.

Note that the Fourier coefficients of different time series correspond to different frequencies, due to the dependency of the number of samples. In order to compare the different time series of \cref{fig:example time series}, linear interpolation is used to retrieve the Fourier coefficients at $\unit[0]{Hz}$, $\unit[0.5]{Hz}$, $\unit[1]{Hz}$, and $\unit[1.5]{Hz}$. Furthermore, the resulting coefficients are divided over the number of samples, see \cref{tab:fourier coefficients}.

\begin{table}
	\centering
	\caption{Interpolated Fourier coefficients for the different time series (see \cref{fig:example time series}) divided over the number of samples.}
	\label{tab:fourier coefficients}
	\begin{tabular}{crrr}
    \toprule
    Time    & \multicolumn{3}{c}{Fourier coefficients} \\
    series  & $\unit[1]{Hz}$ & $\unit[2]{Hz}$ & $\unit[3]{Hz}$ \\\otoprule
    $\profile{1}$ & $0.51-1.79\imagi$ & $0.55-0.78\imagi$ & $0.56-0.41\imagi$ \\
    $\profile{2}$ & $0.28-0.93\imagi$ & $0.29-0.42\imagi$ & $0.29-0.22\imagi$ \\
    $\profile{3}$ & $0.26-0.88\imagi$ & $0.28-0.40\imagi$ & $0.28-0.21\imagi$ \\
    $\profile{4}$ & $0.24-1.01\imagi$ & $0.29-0.44\imagi$ & $0.30-0.23\imagi$ \\
    \bottomrule
\end{tabular}

\end{table}

\Cref{tab:dissimilarities} shows resulting feature-based distance of \cref{eq:eucl features} using the Fourier coefficients of \cref{tab:fourier coefficients} and weights 
$w_{0}=\unit[4.5 \cdot 10^{-2}]{s^2/m^2}$, $w_{1}=\unit[1.1 \cdot 10^{1}]{s^2/m^2}$, $w_{2}=\unit[2.1 \cdot 10^{2}]{s^2/m^2}$, and $w_{3}=\unit[4.9 \cdot 10^{2}]{s^2/m^2}$. 
The weights are the reciprocal of variance of the different coefficients. The variation is based on all time series (around 3000). 
Based on this, it can be seen that $\profile{1}$ differs quite a bit from the other time series. The time series $\profile{2}$ and $\profile{3}$ show the lowest dissimilarity, which matches the (subjective) observation when examining \cref{fig:example time series}. However, although $\profile{4}$ differs significantly from the other time series, it also gets a relatively low dissimilarity score. This shows the disadvantage of this method in that it does not take the differences of the time series' length into account.

\begin{table}
	\centering
	\caption{Dissimilarity measures, according to \cref{eq:eucl features}, with Fourier coefficients (see \cref{sec:fourier coefficient}), custom scenario features (see \cref{sec:custom features}), parameters of an auto-regression model fit (see \cref{sec:auto regression models}), and splines parameters (see \cref{sec:spline models}).}
	\label{tab:dissimilarities}
	\begin{tabular}{ccrrr}
    \toprule
    \multicolumn{2}{c}{Time series} & \multicolumn{3}{c}{Features} \\
    \textbf{x} & \textbf{y} & Fourier coefficients & Auto-regression coefficients & Custom features \\ \otoprule
    $\profile{1}$ & $\profile{2}$ & 1.0e+00 & 2.7e-01 & 3.4e+00 \\
    $\profile{1}$ & $\profile{3}$ & 1.1e+00 & 2.0e-01 & 4.0e+00 \\
    $\profile{1}$ & $\profile{4}$ & 9.9e-01 & 4.7e-01 & 1.6e+01 \\
    $\profile{2}$ & $\profile{3}$ & 6.0e-02 & 3.9e-01 & 1.0e+00 \\
    $\profile{2}$ & $\profile{4}$ & 8.8e-02 & 2.0e-01 & 1.6e+01 \\
    $\profile{3}$ & $\profile{4}$ & 1.3e-01 & 5.9e-01 & 1.7e+01 \\
    \bottomrule
\end{tabular}

\end{table}


\subsubsection{Custom scenario features}
\label{sec:custom features}

Instead of using the Fourier coefficients, it is possible to use scenario parameters. Using custom parameters provides a freedom in choosing parameters that matter. A disadvantage is that these parameters need to be determined for every scenario class, as the parameters are specific for a certain scenario class. For example, the total reduced speed is only relevant when an braking activity is involved. In \cref{tab:custom parameters}, four parameters are listed for describing the time series of \cref{fig:example time series}, namely the starting speed $\parvstart$, the final speed $\parvend$, the reduced speed $\pardv=\parvstart-\parvend$, and the total time $\pardt$ of the time series.

\begin{table}
	\centering
	\caption{Custom parameters for describing the time series of \cref{fig:example time series}.}
	\label{tab:custom parameters}
	\begin{tabular}{crrrr}
    \toprule
    Time & \multicolumn{4}{c}{Parameters} \\
    series & $\parvstart\,[\unit{m/s}]$ & $\parvend\,[\unit{m/s}]$ & $\pardv\,[\unit{m/s}]$ & $\pardt\,[\unit{s}]$ \\\otoprule
    $\profile{1}$ & 13.5 & 0.0 & 13.5 & 4.3 \\
    $\profile{2}$ & 13.1 & 0.0 & 13.1 & 7.6 \\
    $\profile{3}$ & 13.5 & 0.0 & 13.5 & 8.3 \\
    $\profile{4}$ & 8.5 & 4.5 & 3.9 & 2.4 \\
    \bottomrule
\end{tabular}

\end{table}

For comparing the different time series, the parameters $\parvend$, $\pardv$, and $\pardt$ are used. The associated weights are $w_{0}=9.3 \cdot 10^{-2}$, $w_{1}=5.2 \cdot 10^{-2}$, and $w_{2}=3.3 \cdot 10^{-2}$, respectively.
The weights are the reciprocal of the variance of the different parameters based on all time series. The resulting dissimilarity measure is shown in \cref{tab:dissimilarities}. Note that using the aforementioned parameters, only the starting speed, end speed, and total time determines the dissimilarity between two time series. In order to take the `shape' of the time series into account, other parameters are needed. One possibility is to combine the custom parameters with the Fourier coefficients. 


\subsection{Model-based measures}
\label{sec:model based}

With model-based measures, time series models are employed. The idea is to compare the parameters obtained from learning a model of the two time series. For the comparison, we will use the Euclidean distance of \cref{eq:eucl features} \cite{serra2014empirical}. It is shown that this method works well for auto-regression (AR) models \cite{piccolo1990distance, corduas2008autoregressive}, which is demonstrated in the next section. The use of spline models is explained in \cref{sec:spline models}.


\subsubsection{Auto-regression models}
\label{sec:auto regression models}
An AR model approximates the time series using the model 
\begin{equation}
	\label{eq:auto-regression}
	\hat{x}[i] = a_0 + \sum_{j=1}^{\eta-1} a_j \hat{x}[i-j],
\end{equation}
where $\eta$ denotes the number of parameters. \Cref{tab:ar coefficients} lists the parameters of the AR model with $\eta=3$ for the time series of \cref{fig:example time series}. The parameters are computed using ordinary least squares (an alternative would be to use the Yule-Walker equations, see \cite{gardner1988spectral}). As with the previous measures, the weights are the reciprocal of the variance of the parameters, based on the model fits of all available time series. The weights are $w_{0}=\unit[5.4 \cdot 10^{1}]{s^2/m^2}$, $w_{1}=7.4$, and $w_{2}=7.1$, respectively.

\begin{table}
	\centering
	\caption{Coefficients of auto-regression model \cref{eq:auto-regression} with $\eta=3$ for the three time series of \cref{fig:example time series}.}
	\label{tab:ar coefficients}
	\begin{tabular}{crrr}
    \toprule
    Time    & \multicolumn{3}{c}{Auto-regression coefficients} \\
    series & $a_0\,[\unit{m/s}]$ & $a_1$ & $a_2$ \\ \otoprule
    $\profile{1}$ & -0.04 & 1.50 & -0.50 \\
    $\profile{2}$ & -0.03 & 1.58 & -0.57 \\
    $\profile{3}$ & -0.05 & 1.21 & -0.21 \\
    $\profile{4}$ & 0.03 & 1.98 & -1.00 \\
    \bottomrule
\end{tabular}

\end{table}

\Cref{tab:dissimilarities} shows the dissimilarity measures based on the AR parameters. It does not match with the expectation that the dissimilarity score between the time series $\profile{2}$ and $\profile{3}$ would be low. For example, the dissimilarity score between the time series $\profile{1}$ and $\profile{2}$ is lower. Hence, the AR parameters seem to be inappropriate for measuring the dissimilarity between two time series for this application.


\subsubsection{Spline models}
\label{sec:spline models}

Splines can be employed to model and/or smooth time series data \cite{dierckx1993splines, deboor1978practical}. In its most general definition, a spline is a piece-wise polynomial function. It is defined by a degree $d$, which is the maximal degree of the polynomial functions, a set of contiguous intervals with knots $k_0 < k_1 < \ldots < k_n$ as start and end points, and polynomials on each of these intervals. In addition to the basic elements, continuity conditions can be set at the points where the polynomials are stitched together. 

For modeling the time series of the velocity, B-splines are used. B-splines are splines written as a linear combination of $n+d$ basis functions, such that the spline is continuous up to the $(d-1)$-th derivative. Efficient algorithms for the computation of the optimal representation of a curve by B-splines are described in \cite{dierckx1993splines}. We use $d=3$ and $n=2$. The knots are evenly distributed.

The spline parameters, denoted by $c_1$ through $c_5$, of the time series of \cref{fig:example time series} are shown in \cref{tab:spline coefficients}. Because the time series have different length, the knots are placed at different locations for each time series. As a result, the spline coefficients are invariant regarding the length of the time series. Therefore, the length of the time series is added as a parameter when using the euclidean distance of \cref{eq:eucl features} between the parameters. For completeness, the length of the time series is shown in \cref{tab:spline coefficients}. 
The associated weight for the length of the time series is $w_0=\unit[9.3 \cdot 10^{-2}]{s^{-2}}$.
For the spline coefficients $c_1$ through $c_{5}$, the weights are $w_{1}=\unit[4.0 \cdot 10^{-2}]{s^2/m^2}$, $w_{2}=\unit[4.0 \cdot 10^{-2}]{s^2/m^2}$, $w_{3}=\unit[4.2 \cdot 10^{-2}]{s^2/m^2}$, $w_{4}=\unit[3.9 \cdot 10^{-2}]{s^2/m^2}$, and $w_{5}=\unit[3.4 \cdot 10^{-2}]{s^2/m^2}$, respectively.

\begin{table}
	\centering
	\caption{Coefficients of the B-splines of the time series of \cref{fig:example time series}. Because the length of the time series is not reflected by the spline coefficients, the duration of the time series is added to the table.}
	\label{tab:spline coefficients}
	\begin{tabular}{crrrrrr}
    \toprule
    Time & $\pardt\,[\unit{s}]$ & \multicolumn{5}{c}{Spline parameters} \\
    series & & $c_{1}\,[\unit{m/s}]$ & $c_{2}\,[\unit{m/s}]$ & $c_{3}\,[\unit{m/s}]$ & $c_{4}\,[\unit{m/s}]$ & $c_{5}\,[\unit{m/s}]$ \\\otoprule
    $\profile{1}$ & 4.3 & 13.7 & 12.4 & 7.2 & 1.5 & 0.2 \\
    $\profile{2}$ & 7.6 & 13.1 & 12.9 & 5.7 & 3.2 & 0.1 \\
    $\profile{3}$ & 8.3 & 13.5 & 13.0 & 6.7 & 3.9 & 0.1 \\
    $\profile{4}$ & 2.4 & 8.5 & 8.4 & 6.8 & 5.0 & 4.5 \\
    \bottomrule
\end{tabular}

\end{table}

The resulting dissimilarity measures using the length of the time series and the spline coefficients, as shown in \cref{tab:spline coefficients}, is shown in \cref{tab:dissimilarities}. The dissimilarity score between $\profile{2}$ and $\profile{3}$ is, as expected, low.


\subsection{Elastic measures}

Elastic measures allows an elastic shifting of the time axis, to accommodate sequences that are similar but out of phase \cite{keogh2005exactDTW}. Elastic measures have extensively been used for speech recognition \cite{rabiner1993speech, muller2007information}. Due to its success, it is nowadays used in many other applications \cite{serra2014empirical, wang2013distancemeasures}. Dynamic Time Warping (DTW) \cite{sakoe1978dynamic, berndt1994using} belongs to the group of elastic dissimilarity measures \cite{wang2013distancemeasures}.

With DTW, a cost function $D(x[i], y[j])$ is used to compare samples $x[i]$ and $y[j]$, with $i \in \{0, \ldots, N-1\}$ and $j \in \{0, \ldots, M-1\}$. Usually, the cost function is simply the square difference, i.e. $D(x[i], y[j]) = (x[i]-y[j])^2$, or the absolute difference, i.e., $D(x[i], y[j])=|x[i] - y[j]|$ \cite{berndt1994using}. Next, a warping path $(\textbf{w}_x, \textbf{w}_y)$, with elements $(w_x[k], w_y[k]),\,k \in \{0, \ldots, K-1\}$, is defined, such that the dissimilarity measure between two time series is the accumulated sum of the cost function along the optimal warping path, i.e.,
\begin{equation} \label{eq:dtw cost}
	\ddtw(\textbf{x}, \textbf{y}) = \min_{\textbf{w}_x, \textbf{w}_y} \sum_{k=0}^{K-1} D(x[w_x[k]], y[w_y[k]]).
\end{equation}

The warping functions are typically subjected to the following constraints \cite{keogh2005exactDTW}:
\begin{itemize}
	\item \emph{Boundary conditions:} The warping functions are required to start at the first sample of both time series and end at the last sample of both time series, i.e. $w_x[0]=w_y[0]=0$, $w_x[K-1]=N-1$, and $w_y[K-1]=M-1$. 
	\item \emph{Continuity:} To restrict the allowable steps in the warping path to adjacent cells (including diagonally adjacent cells), the constraints $w_x[k] - w_x[k-1] \leq 1$ and $w_y[k] - w_y[k-1] \leq 1$ are used with $k=1, \ldots, K-1$.
	\item \emph{Monotonicity:} The points of the warping path are forces to be monotonically spaced in time using the constraints $w_x[k] - w_x[k-1] \geq 0$ and $w_y[k] - w_y[k-1] \geq 0$ for $k=1, \ldots, K-1$. 
\end{itemize}

To prevent pathological warping paths, the warping path is often constrained in a global sense by limiting how far it may stretch from the diagonal $w_x[k]=w_y[k]$ \cite{keogh2005exactDTW}. The Sakoe and Chiba band \cite{sakoe1978dynamic} and Itakura parallelogram \cite{itakura1975minimum} are most frequently used. Other ways to limit the warping is to add an extra cost for the total stretch from the diagonal, i.e., a cost as a function of $|w_x[k] - w_y[k]|$ to the minimization of \cref{eq:dtw cost}, cf., \cite{roberts1987enhancements, wang1997alignment}.

\Cref{tab:dissimilarities} shows the total cost computed using \cref{eq:dtw cost} divided over $N+M$. It shows a low score for $\ddtw(\profile{2}, \profile{3})$, which is expected because $\profile{2}$ and $\profile{3}$ look very similar (see \cref{fig:example time series}). Furthermore, $\ddtw(\profile{1}, \profile{2})$ and $\ddtw(\profile{1}, \profile{3})$ have low scores. This is because the start and end speeds are similar. Therefore, when the warping is applied, the time series look very similar, see \cref{fig:example dtw}.

\begin{figure}
	\centering
	\setlength\figureheight{0.4\linewidth}
	\setlength\figurewidth{0.7\linewidth}
	% This file was created by matplotlib2tikz v0.6.14.
\begin{tikzpicture}

\definecolor{color1}{rgb}{1,0.498039215686275,0.0549019607843137}
\definecolor{color0}{rgb}{0.12156862745098,0.466666666666667,0.705882352941177}

\begin{axis}[
xlabel={$k$},
ylabel={Speed [km/h]},
xmin=-4.15, xmax=87.15,
ymin=-0.676, ymax=14.196,
width=\figurewidth,
height=\figureheight,
tick align=outside,
tick pos=left,
xmajorgrids,
x grid style={lightgray!92.026143790849673!black},
ymajorgrids,
y grid style={lightgray!92.026143790849673!black},
legend style={draw=white!80.0!black},
legend cell align={left},
legend entries={{$\textbf{v}_1$},{$\textbf{v}_3$}}
]
\addlegendimage{no markers, color0}
\addlegendimage{no markers, color1}
\addplot [semithick, color0, mark=*, mark size=1, mark options={solid}, only marks]
table {%
0 13.52
1 13.5099977498594
2 13.5099977498594
3 13.3099966247891
4 13.3099966247891
5 13.3099966247891
6 13.1599895783992
7 13.1599895783992
8 12.9099906047691
9 12.9099906047691
10 12.649959240875
11 12.649959240875
12 12.3799513653828
13 12.3799513653828
14 12.0599346281397
15 12.0599346281397
16 12.0599346281397
17 11.7499154429195
18 11.7499154429195
19 11.4699158368463
20 11.4699158368463
21 11.1099068362838
22 11.1099068362838
23 10.7949549581562
24 10.7949549581562
25 10.4299515232865
26 10.4299515232865
27 10.1299066003482
28 10.1299066003482
29 9.76989803410774
30 9.76989803410774
31 9.44486661995407
32 9.12984864843619
33 9.12984864843619
34 9.12984864843619
35 8.72493720331242
36 8.32991331122718
37 8.32991331122718
38 7.98480630039439
39 7.98480630039439
40 7.98480630039439
41 7.62981712672889
42 7.62981712672889
43 7.28483020320402
44 7.28483020320402
45 6.92979800053355
46 6.92979800053355
47 6.92979800053355
48 6.51990004737112
49 6.51990004737112
50 6.20973788094142
51 6.20973788094142
52 6.20973788094142
53 5.72975396488627
54 5.72975396488627
55 5.72975396488627
56 5.35981499501638
57 5.35981499501638
58 4.96973282804373
59 4.96973282804373
60 4.5897063595452
61 4.5897063595452
62 4.5897063595452
63 4.19466363358235
64 4.19466363358235
65 3.81966622914132
66 3.81966622914132
67 3.41487309404241
68 3.41487309404241
69 3.41487309404241
70 3.029868466762
71 2.72972905543558
72 2.72972905543558
73 2.43974109566171
74 2.17971330445119
75 1.90970542895899
76 1.90970542895899
77 1.61484808590058
78 1.32979961905576
79 1.06470324790128
80 0.829719324563913
81 0.609807593238133
82 0.418966421631237
83 0
};
\addplot [semithick, color1, mark=*, mark size=1, mark options={solid}, only marks]
table {%
0 13.52
1 13.5024986898786
2 13.4599984012158
3 13.3749877623932
4 13.339986814965
5 13.2549875650122
6 13.169975031088
7 13.0799882162525
8 12.9599456808151
9 12.8199847227294
10 12.6799575631365
11 12.5599439833408
12 12.449928705367
13 12.3299113356886
14 12.2099400949398
15 12.1199744389397
16 11.9699583719875
17 11.8398822031638
18 11.7198902957744
19 11.5699011704795
20 11.3898610439473
21 11.259927127024
22 11.0899237313226
23 10.9148207585339
24 10.7497923533932
25 10.5597400469769
26 10.3998420756563
27 10.2147452236211
28 10.0198336458553
29 9.79482735921659
30 9.60982176341755
31 9.38481547677722
32 9.2198091920216
33 9.02977514384154
34 8.95961715650411
35 8.67975721927255
36 8.46984377312979
37 8.28974324711024
38 8.12973645461187
39 7.94979725048586
40 7.81975751534602
41 7.63964461243898
42 7.47981781756124
43 7.33977628864905
44 7.18988559811299
45 7.05984406920148
46 6.89972117718984
47 6.74463343070855
48 6.60970860678816
49 6.44974514196542
50 6.27995667492372
51 6.14469078330582
52 6.01972630914031
53 5.86962789412376
54 5.72966812399489
55 5.5647587020236
56 5.39980356473154
57 5.24960034344291
58 5.09964366851783
59 4.9296897736018
60 4.76952658225369
61 4.59473274645391
62 4.40972850008857
63 4.25977920770587
64 4.07972001263189
65 3.92449063474453
66 3.74959959339546
67 3.60470975159601
68 3.42952880800507
69 3.24976101137892
70 3.05975761403809
71 2.86950804334438
72 2.67931368736811
73 2.46430435125915
74 2.22942319421175
75 1.99961037038475
76 1.76960527406612
77 1.53426645678292
78 1.30928698175165
79 1.08427764564269
80 0.889268309531239
81 0.659730535789912
82 0.519454276296623
83 0
};
\end{axis}

\end{tikzpicture}
	\caption{The time series $\profile{1}$ and $\profile{3}$ look very similar after warping using Dynamic Time Warping.}
	\label{fig:example dtw}
\end{figure}


\printbibliography

\end{document} 